\documentclass[12pt]{article}
\usepackage[utf8]{inputenc}
\usepackage[T1]{fontenc}
\usepackage{hyperref}
\title{Puri Language Resources for K-12 Curriculum Development}
\author{Compiled by Claude Sonnet 5.5 (claude-sonnet-5-5) for Andr\'as Kornai}
\date{}
\begin{document}
\maketitle

\section{Introduction}
Puri (ISO 639-3 prr; Purian) is a dormant language of eastern Brazil. It is attested only through a few 19th-century word lists. One source notes that Puri people have reconstructed the language and that some second-language speakers exist. The search found no teaching materials of any kind.

\section{Primary Educational Resources}
None found.

\section{Linguistic and Reference Materials}
\begin{itemize}
\item Wikipedia, Puri language: \url{https://en.wikipedia.org/wiki/Puri_language}
\item Wikipedia, Purian languages (Puri, Coroado, and the disputed Corop\'o): \url{https://en.wikipedia.org/wiki/Purian_languages}
\item Silva Neto (2007), a University of Bras\'ilia thesis reconstructing Proto-Purian: \url{https://repositorio.unb.br/bitstream/10482/2328/1/2007_AmbrosioPereiradaSilvaNeto.pdf}
\item Ramirez, Vegini and Fran\c{c}a (2015), LIAMES, reviewing the Corop\'o and Puri word-list data: \url{https://periodicos.sbu.unicamp.br/ojs/index.php/liames/article/view/8642302}
\end{itemize}

\section{Religious and Community Resources}
None found.

\section{Digital and Audio Resources}
None found.

\section{Curriculum Development Context}
The documentary base is a handful of early 19th-century word lists (for example Martius 1818 and Eschwege 1815). Any heritage curriculum would be a small revitalisation programme built with the Puri community, with cultural and historical content in Portuguese rather than a language-medium curriculum. Corop\'o is classified by recent work as a separate, Macro-J\^e language, so its lists should not be mixed in.
\end{document}